\documentclass{book}
\usepackage[utf8]{inputenc}
\usepackage[T1]{fontenc}
\usepackage{geometry}
\usepackage{graphicx}
\usepackage{float}
\usepackage{hyperref}
\usepackage[english,indonesian]{babel}

\title{Optimasi Pemrosesan Template LaTeX Menggunakan Pendekatan Berbasis AI}
\author{Academic Support Agent}
\date{\today}

\begin{document}
\maketitle

\begin{abstract}

\end{abstract}

\chapter{PENDAHULUAN}


\chapter{TINJAUAN PUSTAKA}


% Original content (commented out):
% % PLANTILLA LaTeX VIU - mayo 2020
% Basado en plantilla original de Rafael Padilla (2011)
% Modificaciones de Juande Santander-Vela (2020)
%----------------------------------------------------------
% PREAMBULO:
\documentclass[a4paper,11pt]{book}
%----------------------------------------------------------
% PAQUETES:
% Language

% In moder computers, this should be set to UTF8
\usepackage[utf8]{inputenc}
%\usepackage[latin1]{inputenc} // only for non-UTF supporting 
\usepackage[spanish]{babel}

% SECTIONS WITH SmallCaps
\usepackage[T1]{fontenc}

% PAQUETE PARA AÑADIR URLs
\usepackage{url}
%% Define a new 'leo' style for the package that will use a smaller font.
\makeatletter
\def\url@leostyle{%
  \@ifundefined{selectfont}{\def\UrlFont{\sf}}{\def\UrlFont{\small\ttfamily}}}
\makeatother
%% Now actually use the newly defined style.
\urlstyle{leo}

% PARA HACER CLICKABLES LAS URLs EN EL PDF
\usepackage{hyperref}
% SETUP COLORES PDF
\hypersetup{
    colorlinks = false,
    linkcolor = red,
    anchorcolor = red,
    citecolor = blue,
    filecolor = red,
    urlcolor = red
}

% Formulas y simbolos de matematicas:
\usepackage{amsmath}
\usepackage{amssymb}
\usepackage{gensymb}
\usepackage{wasysym}
\usepackage{mathrsfs}

% Graficos
\usepackage{graphicx}
\usepackage{float}
\usepackage{lscape}
\usepackage{wrapfig}

% NOTA:
% Si compilamos en DVI->dvips->ps2pdf tendremos que usar: EPS o PS (tamin BMP)
% Si compilamos directamente en pdf: JPEG o PNG

% INDEX
\usepackage{makeidx}

% MOSTRAR BIBLIOGRAFIA EN INDICE
\usepackage[nottoc]{tocbibind}

% PAQUETE PARA LA BIBLIOGRAFIA: 
% BibTeX
\usepackage[sectionbib]{natbib}  % Cross-reference package (Natural BiB)
\bibpunct{(}{)}{;}{a}{}{,}
%\usepackage{chapterbib}  % Put References at the end of each chapter
% Archivo con la bibliografia: ./files/TFM/Biblio/Bib.bib
% Compilar: Normal, BibTeX y 2 veces mas como Normal

% Tablas
\usepackage{multirow} % Para unir varias filas
\usepackage{longtable} % Para crear tablas de mas de una pagina

%----------------------------------------------------------
% Encabezados y pies de pagina
\usepackage{fancyhdr}
% Formato para los encabezados y pies de pagina
% En lo siguiente, fancyhead sirve para configurar la cabecera, fancyfoot para el pie.
% Justificación: C=centered, R=right, L=left, (nada)=LRC
% Página: O=odd, E=even, (nada)=OE

% Formato FANCY (Usando \thispagestyle{fancy})
\pagestyle{fancy}
% Borrar todos los ajustes previos
\fancyhf{}
% Redefinimos el comando \chaptermark (mostrar Capitulo en roman)  -> Aparece en \leftmark
%\renewcommand{\chaptermark}[1]{\markboth{#1}{}}
% Redefinimos el comando \sectionmark (mostrar Seccion en roman) -> Aparece en \rightmark
%\renewcommand{\sectionmark}[1]{\markright{#1}}
% Encabezado
\fancyhead[LE]{Valencian International University / VIU}
\fancyhead[RO]{Máster en Astronomía y Astrofísica}
% Pie de pagina
\fancyfoot[LE, RO]{\thepage}
% Tamaño de las lineas del encabezado y pie de pagina
\renewcommand{\headrulewidth}{0.5pt}
\renewcommand{\footrulewidth}{0.5pt}
% Si quitamos las lineas, mejor poner el texto de los encabezados y pie de pagina en negrita o cursiva
% Formato PLAIN (Usando \thispagestyle{plain}) CREO QUE ES EL QUE USAN LOS CHAPTERS...
\fancypagestyle{plain}{
\fancyhf{} % borrar todos los ajustes
\renewcommand{\headrulewidth}{0pt}
\renewcommand{\footrulewidth}{0.5pt}
\fancyfoot[LE, RO]{\thepage}
}
% Formato EMPTY (Usando \thispagestyle{empty})
%----------------------------------------------------------


%Paquete Theorem:
\usepackage{theorem}
%\usepackage{amsthm}
%Defs. del TEOREM:      %%%%%%%%%%%%%%%%%%%%%%%%%%%%%%%%%%%%%%
\theoremstyle{break}
\newtheorem{Def}{Definición}
\newtheorem{Prob}{Problema}
\newtheorem{Prop}[Def]{Propiedad}
\newtheorem{Propo}{Proposición}
\newtheorem{Cor}{Corolario}
\newtheorem{Teo}{Teorema}
\newtheorem{Alg}{Algoritmo}

\newcommand{\ind}{\textrm{ind}}
\newcommand{\Ind}{\textrm{Ind}}
\newcommand{\tab}{\hspace{5mm}}

%\newtheorem{Ejer}{Ejercicio}
%\newtheorem{Ejem}{Ejemplo}


{\theorembodyfont{\rmfamily} \newtheorem{Ejer}{Ejercicio}}
{\theorembodyfont{\rmfamily} \newtheorem{Ejem}{Ejemplo}}

\newenvironment{demo}{{\bf
Demostración}}{\begin{flushright}Q.E.D.\end{flushright}}
\newenvironment{nota}{{\bf Nota: }}{\vspace{12pt}}
%%%%%%%%%%%%%%%%%%%%%%%%%%%%%%%%%%%%%%%%%%%%%%%%%

% PAQUETE PARA ESPACIADOS
\usepackage{setspace}

% PARA CORREGIR
\usepackage[modulo,switch]{lineno} 

% PAQUETE PARA HACER CAPTIONS MAS VISUALES
\usepackage[font=small,format=plain,labelfont=bf,up,textfont=it,up]{caption}

% DEFINICIONES
\def\d{$\,\rm day$ }
%\def\h{$\,\rm h$ }
%\def\min{$\,\rm min$ }
%\def\seg{$\,\rm s$ }
\def\km{$\,\rm km$ }
\def\m{$\,\rm m$ }
\def\mm{$\,\rm mm$ }
\def\ppm{$\,\rm ppm$ }
\def\kg{$\,\rm Kg$ }
\def\K{$\,\rm K$ }
\def\arcsec{$\,\rm arcsec$ }
\def\pixels{$\,\rm pixels$ }
\newcommand{\cd}{{\, \rm c\, d^{-1}}}
\newcommand{\h}{{\rm ^{h}}}
\newcommand{\minutos}{{\rm ^{m}}}
\newcommand{\segundos}{{\rm ^{s}}}
\newcommand{\MHz}{\,\rm{MHz}}
\newcommand{\rmdiv}{{\rm div}}
\def\kms{$\, \rm km\, s^{-1}$}
% Comando ION
\newcommand\ion[2]{#1$\;${\scshape{#2}}}%                       % ion, i.e., CII = \ion{C}{ii}

% REVISTAS
\def\aap{Astronomy \& Astrophysics}
\def\aaps{Astronomy \& Astrophysics Supplement Series}
\def\apj{Astrophysical Journal}
\def\apjl{Astrophysical Journal Letters}
\def\apjs{Astrophysical Journal Supplement Series}
\def\mnras{Monthly Notices of the Royal Astronomical Society}
\def\nat{Nature}
\def\araa{Annual Review of Astronomy and Astrophysics}
\def\pasj{Publications of the Astronomical Society of Japan}
\def\araa{Annual Review of Astronomy \& Astrophysics}
\def\aj{Astronomical Journal}
\def\apss{Astrophysics and Space Science}
\def\pasp{Publications of the Astronomical Society of the Pacific}



% MAKES
\makeindex

% MAIN
\begin{document}
% PARA CORREGIR
%\linenumbers
% Separacion entre parrafos
\parskip 15pt

% Separacion entre lineas
\onehalfspacing
% \singlespacing
% \doublespacing



%%%%%%%%%%%%%%%%%%%%%%%%%%%%%


% TITULO DEL TRABAJO
\thispagestyle{empty}

\input{files/portada.tex}

\newpage


% INDICE
\pagestyle{fancy}
\pagenumbering{roman}
\tableofcontents


% CAPITULOS
\newpage
\pagestyle{fancy}
\pagenumbering{arabic}

% CAPITULO 1
\chapter{Título Capítulo 1}
\label{chapter:01}
\input{./files/cap01.tex}

% CAPITULO 2
\chapter{Título Capítulo 2}
\label{chapter:02}
\input{./files/cap02.tex}


% BIBLIOGRAFIA
\bibliographystyle{./sty/astron}
\bibliography{./biblio/tfm_biblio}

%%%%%%%%%%%%%%%%%%%%%%%%%%%%%

\end{document}


\end{document}